% ============================================================================
% Chapitre 2 --- Spécifications fonctionnelles et techniques
% ============================================================================
\chapter{Spécifications fonctionnelles et techniques}
\label{chap:specs}

\section{Spécifications fonctionnelles}

Le cahier des charges propose une liste de fonctionnalités possibles. Le tableau~\ref{tab:fonctionnel} indique comment chacune est couverte par le projet.

\begin{table}[H]
\centering
\caption{Couverture des fonctionnalités du cahier des charges.}
\label{tab:fonctionnel}
\begin{tabularx}{\textwidth}{@{}p{4.2cm} X c@{}}
\toprule
\textbf{Fonctionnalité} & \textbf{Réalisation dans le projet} & \textbf{Statut} \\
\midrule
Automatisation des processus & Enchaînement automatique des postes, transfert par convoyeur, blocage et reprise sans intervention & Réalisé \\
Collecte et analyse de données en temps réel & État complet diffusé 10~fois/s, \ac{KPI} calculés en continu, historique SQLite par minute simulée & Réalisé \\
Interconnexion des équipements et systèmes & Simulation, backend, scène \ac{3D} et tableau de bord reliés par un protocole \ac{WS} unique & Réalisé \\
Interface homme-machine intuitive & Scène \ac{3D} navigable, voyants d'état, tableau de bord, onglet d'aide, manuel utilisateur & Réalisé \\
Maintenance prédictive & Pannes simulées, \ac{MTBF} et \ac{MTTR} observés par poste~; pas de modèle prédictif & Partiel \\
\bottomrule
\end{tabularx}
\end{table}

Les fonctions détaillées attendues par l'utilisateur sont les suivantes~:

\begin{itemize}
  \item \textbf{F1 --- Visualiser} la ligne en \ac{3D}~: machines, convoyeur, produits en mouvement, état de chaque machine (en marche, en attente, bloquée, en panne), avancement du produit en cours.
  \item \textbf{F2 --- Superviser}~: cadence réelle et théorique, produits finis, rejets, en-cours, goulot, état et statistiques de chaque poste.
  \item \textbf{F3 --- Piloter}~: pause, vitesse de simulation ($\times$1 à $\times$20), remise à zéro, modification des temps de cycle, nombre de machines en parallèle, convoyeur adaptatif, activation des pannes.
  \item \textbf{F4 --- Simuler des incidents}~: provoquer ou réparer une panne sur un poste choisi.
  \item \textbf{F5 --- Comparer}~: simuler une heure de chaque scénario dans des conditions identiques et afficher les écarts par rapport à la situation initiale.
  \item \textbf{F6 --- Historiser}~: conserver chaque exécution, consulter ses courbes et événements, exporter les données en \ac{CSV}.
\end{itemize}

\section{Spécifications techniques}

\subsection{Technologies retenues}

\begin{table}[H]
\centering
\caption{Choix technologiques.}
\label{tab:techno}
\begin{tabularx}{\textwidth}{@{}l l X@{}}
\toprule
\textbf{Brique} & \textbf{Technologie} & \textbf{Justification} \\
\midrule
Simulation & Python 3.12 & Langage maîtrisé par toute l'équipe, lisible pour un modèle métier, sans dépendance \\
Backend & FastAPI + Uvicorn~\cite{fastapi} & Serveur asynchrone léger qui gère à la fois \ac{WS} et \ac{REST} \\
Communication & WebSocket~\cite{rfc6455} & Canal bidirectionnel persistant et à faible latence~: l'état descend, les commandes remontent \\
Stockage & SQLite & Base \ac{SQL} embarquée dans la bibliothèque standard, un seul fichier, aucun serveur \\
\ac{3D} & Unity~6 (URP), export WebGL~\cite{unity6} & Moteur de jeu standard de l'industrie pour les jumeaux numériques, rendu dans le navigateur sans installation côté poste opérateur \\
Tableau de bord & HTML, CSS, JavaScript & Même page que la scène Unity, relié au même backend~; le tableau de bord et la \ac{3D} reçoivent le même état \\
Validation & pytest & Tests automatisés reproductibles du modèle et de l'\ac{API} \\
\bottomrule
\end{tabularx}
\end{table}

\begin{callout}
\textbf{Pourquoi Unity avec export WebGL~?} Le plan de travail prévoyait Unity, moteur de référence pour les jumeaux numériques. Une première version de la scène avait été écrite avec Three.js~; nous l'avons reconstruite dans Unity pour respecter la consigne et obtenir un rendu plus réaliste (éclairage, matériaux, brouillard, ombres). Pour ne pas séparer la \ac{3D} du tableau de bord, la scène est exportée en \textbf{WebGL} et affichée dans la même page web que le tableau de bord~: le navigateur sert d'unique poste de supervision, sans installation. Un petit pont (\code{LineBridge}) relie la page et la scène (sélection d'une machine dans les deux sens). La contrepartie est un build plus lourd (environ 12~Mo) et un premier chargement de quelques secondes.
\end{callout}

\subsection{Normes et standards}

Le projet s'appuie sur des standards ouverts~: le protocole WebSocket (RFC~6455~\cite{rfc6455}) pour la communication, \ac{JSON} pour le format des messages, HTTP/\ac{REST} pour l'historique, \ac{CSV} pour l'export. Dans une version industrielle, la collecte des données machines passerait par \ac{OPC} (IEC~62541~\cite{opcua}), standard de fait de l'interopérabilité industrielle~; l'architecture actuelle isole la source de données dans un seul module (\code{sim.py}) précisément pour permettre ce remplacement.

\subsection{Sécurité et cybersécurité}

Le serveur écoute uniquement sur l'interface locale (\code{127.0.0.1})~: il n'est pas exposé au réseau. Les commandes reçues sont validées et bornées côté serveur (numéro de poste, temps de cycle entre 1 et 60~s, 1 à 3 machines, vitesse entre $\times$0{,}25 et $\times$20)~; une commande invalide renvoie un message d'erreur sans modifier l'état. Aucune donnée personnelle n'est traitée. Une mise en réseau réelle demanderait en plus une authentification des opérateurs, un chiffrement (\code{wss://}) et des droits différenciés entre consultation et pilotage~; ces points sont listés dans les perspectives.

\subsection{Performances attendues}

\begin{table}[H]
\centering
\caption{Exigences de performance et valeurs obtenues.}
\label{tab:perf}
\begin{tabularx}{\textwidth}{@{}X p{3.1cm} p{3.6cm}@{}}
\toprule
\textbf{Exigence} & \textbf{Cible} & \textbf{Obtenu} \\
\midrule
Fréquence de mise à jour de l'interface & $\geq$ 5~Hz & 10~Hz \\
Pas de simulation & $\leq$ 0{,}1~s & 0{,}05~s \\
Taille d'un message d'état & $<$ 10~ko & $\approx$ 3~ko \\
Comparaison de 5 scénarios sur 1~h simulée & $<$ 5~s & $\approx$ 0{,}3~s \\
Écart débit simulé / débit théorique en régime & $\leq$ 5~\% & vérifié par test \\
Écart entre la mesure avec la scène Unity et la simulation hors ligne & $\leq$ 2~\% & 1{,}3~\% au maximum \\
Rendu \ac{3D} & fluide (cible 30~images/s ou plus) & à mesurer sur le poste de démonstration \\
\bottomrule
\end{tabularx}
\end{table}
